\begin{abstract}
Instruction-tuning datasets for cybersecurity have, until recently, been
assembled with the help of commercial, cloud-hosted language models. We ask a
narrower and more uncomfortable question: what happens when the entire
construction and fine-tuning loop is moved onto a single consumer
Apple~Silicon machine, with no external API in the path and no provider side
safety filter shaping the corpus? We describe \systemname, a fully on-device
reimplementation of a published cybersecurity data pipeline built on Apple's
MLX framework. Building on the work of
ElZemity et al.~\cite{cyberllminstruct}, we reproduce a defensive instruction
dataset locally and characterize a second, offensive branch that
removes the pipeline's safety-alignment stage and trains models, including
\emph{abliterated} (safety-stripped) bases, to produce exploitation guidance
for vulnerable code. We then measure a full matrix of alignment (aligned vs.\
abliterated) $\times$ fine-tuning (base vs.\ offensively fine-tuned) across
three open model families on capability, safety, and audit quality. The results
are stark: abliteration removes refusals at almost no capability cost, and
offensive fine-tuning drives refusal to zero and attack success rate to
$96$--$100\%$ on \emph{every} base, including those that still refused a fifth
of requests while aligned. The whole procedure runs on a single desktop at no
marginal cost: a complete fine-tune takes hours, not a cloud budget.  We do
\emph{not} release exploitation payloads, offensive weights, the dataset, or the
raw generations: only aggregate metrics. Our contribution is a warning made
concrete and measured: the barrier to a locally hosted, uncensored,
exploitation capable LLM is now a weekend of compute on hardware
that fits under a desk.
\end{abstract}
